\chapter{求导与凸性}\label{ch:g12:deriv}

导数刻画函数的瞬时变化率，它是函数图像
在某点处切线的斜率。本章回顾并推广求导法则，
建立 $f'$ 的符号与 $f$ 的单调性之间的联系，
并引入由二阶导数决定的凸性。

\section{导数}

\begin{definition}[一点处的导数]\label{def:g12:deriv:derivative}
设 $f$ 定义在区间 $I$ 上，$a \in I$。称函数 $f$
在 $a$ 处\emph{可导}\index{导数}\index{可导}，若差商
\[
\frac{f(a+h) - f(a)}{h}
\]
当 $h \to 0$ 时有有限极限。该极限称为
\emph{$f$ 在 $a$ 处的导数}，记作 $f'(a)$。若 $f$ 在 $I$ 的每一点处都可导，
则称函数 $f' \colon x \mapsto f'(x)$ 为
$f$ 的\emph{导函数}。
\end{definition}

\begin{definition}[切线]\label{def:g12:deriv:tangent}
若 $f$ 在 $a$ 处可导，则 $f$ 的
曲线在点 $(a, f(a))$ 处的\emph{切线}\index{切线}为直线，其方程为
\[
y = f(a) + f'(a)(x - a).
\]
\end{definition}

\begin{omfigure}
\begin{tikzpicture}
\begin{axis}[omaxis, clip=false,
  width=9.6cm, height=6.8cm,
  xmin=-0.4, xmax=3.4, ymin=-0.4, ymax=4.3,
  xtick={1,2.5}, xticklabels={$a$,$a+h$}, ytick=\empty]
  \addplot[omProp!40, thick, domain=0.2:2.7] {1 + 1.75*(x-1)};
  \addplot[omProp!70, thick, domain=0.07:2.95] {1 + 1.5*(x-1)};
  \addplot[omProp, thick, domain=-0.12:3.3] {1 + 1.25*(x-1)};
  \addplot[omThm, thick, domain=-0.4:3.3] {x};
  \addplot[omDef, thick, domain=-0.4:2.72] {x^2/2 + 0.5};
  \draw[black, dashed, thin] (axis cs:1,1) -- (axis cs:2.5,1)
    node[midway, below, font=\small] {$h$};
  \draw[black, dashed, thin] (axis cs:2.5,1) -- (axis cs:2.5,3.625)
    node[midway, right, font=\small] {$f(a+h)-f(a)$};
  \fill (axis cs:2.5,3.625) circle (1.5pt);
  \fill (axis cs:2,2.5) circle (1.5pt);
  \fill (axis cs:1.5,1.625) circle (1.5pt);
  \fill (axis cs:1,1) circle (1.5pt)
    node[above left, font=\small] {$A$};
  \node[omThm, font=\small, anchor=west] at (axis cs:3.38,3.3) {切线};
\end{axis}
\end{tikzpicture}

{\small 当 $h$ 减小时，过 $A = (a, f(a))$ 与
$(a+h, f(a+h))$ 的割线（橙色）转向 $A$ 处的切线（红色）：其
斜率 $\frac{f(a+h)-f(a)}{h}$ 趋于 $f'(a)$。}
\end{omfigure}

\begin{proposition}\label{prop:g12:deriv:diffcont}
在 $a$ 处可导的函数在 $a$ 处连续，其逆不成立。
\end{proposition}

\begin{proof}
当 $h \neq 0$ 充分小时，$f(a+h) - f(a) = h \cdot \frac{f(a+h)-f(a)}{h}$。
当 $h \to 0$ 时，右端趋于 $0 \times f'(a) = 0$，故
$f(a+h) \to f(a)$。对于逆命题，$x \mapsto \abs{x}$ 在
$0$ 处连续，但其在 $0$ 处的差商为 $\frac{\abs h}{h} = \pm 1$，
其值取决于 $h$ 的符号，故无极限。
\end{proof}

\subsection{求导法则}

\begin{proposition}[运算法则]\label{prop:g12:deriv:operations}
设 $u, v$ 在 $I$ 上可导，$\lambda \in \R$。则 $u + v$，
$\lambda u$、$uv$ 在 $I$ 上可导，$\frac{u}{v}$ 亦然，
其中 $v \neq 0$，且
\[
(u+v)' = u' + v', \quad (\lambda u)' = \lambda u', \quad
(uv)' = u'v + uv', \quad
\left(\frac{u}{v}\right)' = \frac{u'v - uv'}{v^2}.
\]
\end{proposition}

\begin{proof}[乘积法则的证明]
将 $uv$ 在 $a$ 处的差商写为
\[
\frac{u(a+h)v(a+h) - u(a)v(a)}{h}
= \frac{u(a+h) - u(a)}{h}\,v(a+h) + u(a)\,\frac{v(a+h) - v(a)}{h}.
\]
当 $h \to 0$ 时，由连续性得 $v(a+h) \to v(a)$，
（见\cref{prop:g12:deriv:diffcont}），故右端趋于
$u'(a)v(a) + u(a)v'(a)$。其余法则同理可证。
\end{proof}

\begin{theorem}[链式法则]\label{thm:g12:deriv:chain}
\index{链式法则}
设 $u$ 在 $I$ 上可导且取值于 $J$，$g$ 在
$J$ 上可导。则 $g \circ u$ 在 $I$ 上可导，且
\[
(g \circ u)' = (g' \circ u) \cdot u' .
\]
特别地，对可导的 $u$，有：
\[
(u^n)' = n\,u'\,u^{n-1} \ (n \in \Z), \qquad
\left(\sqrt{u}\right)' = \frac{u'}{2\sqrt{u}} \ (u > 0), \qquad
\left(\eu^{u}\right)' = u'\,\eu^{u}.
\]
\end{theorem}

\begin{proof}[证明概要]
当 $u(a+h) \neq u(a)$ 且 $h$ 充分小时，可写为
\[
\frac{g(u(a+h)) - g(u(a))}{h}
= \frac{g(u(a+h)) - g(u(a))}{u(a+h) - u(a)} \cdot \frac{u(a+h) - u(a)}{h}.
\]
当 $h \to 0$ 时，$u(a+h) \to u(a)$，故第一个因子趋于 $g'(u(a))$，
第二个因子趋于 $u'(a)$。（完整证明须处理
$u(a+h) = u(a)$ 对任意小的 $h$ 成立的情形；此技术性问题
留待大学阶段处理。）
\end{proof}

\begin{example}
以下是在自然定义域上成立的常用导数：
\[
(x^n)' = nx^{n-1}, \quad
\left(\frac{1}{x}\right)' = -\frac{1}{x^2}, \quad
(\sqrt{x})' = \frac{1}{2\sqrt{x}}, \quad
(\cos x)' = -\sin x, \quad
(\sin x)' = \cos x .
\]
第一式由乘积法则用归纳法证明，后两式在
\cref{ch:g12:trigo}中证明。
\end{example}

\section{单调性与极值}

\begin{theorem}[导数的符号与单调性]\label{thm:g12:deriv:variations}
设 $f$ 在区间 $I$ 上可导。
\begin{enumerate}
  \item $f$ 在 $I$ 上递增当且仅当 $f' \geq 0$ 在 $I$ 上。
  \item $f$ 在 $I$ 上递减当且仅当 $f' \leq 0$ 在 $I$ 上。
  \item 若 $f' > 0$ 在 $I$ 上，除有限个导数为零的点外成立，
        则 $f$ 在 $I$ 上严格递增。
\end{enumerate}
\end{theorem}

\begin{proof}[部分证明]
若 $f$ 递增，则每个差商
$\frac{f(a+h)-f(a)}{h}$ 均 $\geq 0$，取极限得
$f'(a) \geq 0$。反向蕴含依赖于中值定理，
该定理留待大学阶段证明，这些结论\emph{在本阶段不加证明地承认}。
\end{proof}

\begin{definition}[局部极值]\label{def:g12:deriv:extremum}
称 $f$ 在 $a$ 处有\emph{局部极大值}\index{极值}，若 $f(x) \leq f(a)$ 对
所有 $x$ 在 $a$ 附近成立；局部极小值可类似地定义。
\end{definition}

\begin{proposition}[一阶条件]\label{prop:g12:deriv:critical}
若 $f$ 在\emph{开}区间 $I$ 上可导，且
在 $a \in I$ 处取得局部极值，则 $f'(a) = 0$。逆命题不成立
（例如 $f(x) = x^3$ 在 $a = 0$ 处）。
\end{proposition}

\begin{proof}
设 $a$ 为局部极大值点。对于充分小的 $h > 0$，
$\frac{f(a+h)-f(a)}{h} \leq 0$，令 $h \to 0^+$ 得 $f'(a) \leq 0$；
当 $h < 0$ 时该商 $\geq 0$，得 $f'(a) \geq 0$。
因此 $f'(a) = 0$。
\end{proof}

\begin{omfigure}
\begin{tikzpicture}
\begin{axis}[omaxis,
  width=8.4cm, height=5.6cm,
  xmin=-2.5, xmax=2.5, ymin=-4.2, ymax=4.2,
  xtick={-1,1}, ytick={-2,2}]
  \addplot[omProp, thick, domain=-1.8:-0.2] {2};
  \addplot[omProp, thick, domain=0.2:1.8] {-2};
  \addplot[omDef, thick, domain=-2.15:2.15] {x^3 - 3*x};
  \fill (axis cs:-1,2) circle (1.5pt);
  \fill (axis cs:1,-2) circle (1.5pt);
\end{axis}
\end{tikzpicture}

{\small 在开区间内部的局部极值点处，切线
（橙色）水平：$f'(a) = 0$。此处 $f(x) = x^3 - 3x$，
在 $-1$ 处有局部极大值，在 $1$ 处有局部极小值。}
\end{omfigure}

\begin{method}[变化表]\label{met:g12:deriv:table}
研究函数 $f$ 时：求其定义域及在端点处的极限；
求 $f'$ 并将其因式分解；确定 $f'$ 在各子区间上的符号；
将所有结果记入变化表并标出极值；
应用双射定理确定方程 $f(x) = k$ 解的个数，
（见\cref{thm:g12:limcont:bijection}）。
\end{method}

\section{凸性}

\begin{definition}[凸函数]\label{def:g12:deriv:convex}
称 $f$ 为区间 $I$ 上的\emph{凸函数}\index{凸性}，若
在 $I$ 上，其图像的任一弦都位于图像上方，即
对所有 $a, b \in I$ 及 $t \in \intcc{0}{1}$，
\[
f\bigl((1-t)a + t b\bigr) \leq (1-t) f(a) + t f(b).
\]
若反向不等式成立（即 $-f$ 为凸函数），则称其为\emph{凹函数}\index{凹函数}。
\end{definition}

\begin{omfigure}
\begin{tikzpicture}
\begin{axis}[omaxis,
  width=8.8cm, height=6cm,
  xmin=-1.9, xmax=2.5, ymin=-0.8, ymax=4.3,
  xtick={-1,1.5}, xticklabels={$a$,$b$}, ytick=\empty]
  \addplot[omThm, thick, domain=-1.25:1.75] {0.5*x + 1.5};
  \addplot[omProp, thick, domain=-0.45:2.3] {x - 0.25};
  \addplot[omDef, thick, domain=-1.6:2.0] {x^2};
  \draw[black, dashed, thin] (axis cs:0.25,0.0625) -- (axis cs:0.25,1.625);
  \fill (axis cs:-1,1) circle (1.5pt);
  \fill (axis cs:1.5,2.25) circle (1.5pt);
  \fill (axis cs:0.5,0.25) circle (1.5pt);
  \node[omThm, font=\small, above left] at (axis cs:1.7,2.4) {弦};
  \node[omProp, font=\small, below right] at (axis cs:1.75,1.5) {切线};
\end{axis}
\end{tikzpicture}

{\small 凸函数：每条弦（红色）都位于图像上方，且
图像位于其各切线（橙色）上方。虚线段
表示 $f\bigl((1-t)a+tb\bigr)$ 与 $(1-t)f(a)+tf(b)$ 之间的差距。}
\end{omfigure}

\begin{theorem}[凸性与导数]\label{thm:g12:deriv:convexity}
设 $f$ 在区间 $I$ 上可导。以下三条等价：
\begin{enumerate}
  \item $f$ 在 $I$ 上为凸函数；
  \item $f'$ 在 $I$ 上递增；
  \item $f$ 的图像位于其每一条切线的上方。
\end{enumerate}
若 $f$ 二阶可导，则上述三条又等价于
$f'' \geq 0$ 在 $I$ 上。
\end{theorem}

\begin{proof}[(2) $\Rightarrow$ (3) 的证明]
取 $a \in I$，并令
$g(x) = f(x) - f(a) - f'(a)(x-a)$，即曲线与
在 $a$ 处切线之间的差距。则 $g$ 可导，且 $g'(x) = f'(x) - f'(a)$。若 $f'$ 递增，
则 $g' \leq 0$ 在 $I \cap \intoc{-\infty}{a}$ 上，且 $g' \geq 0$ 在
$I \cap \intco{a}{+\infty}$ 上：$g$ 在 $a$ 之前递减，之后递增，
故 $g$ 取得最小值 $g(a) = 0$，从而 $g \geq 0$。
其余蕴含在本阶段不加证明地承认；
$f'' \geq 0$ 的等价性由 \cref{thm:g12:deriv:variations} 应用于 $f'$ 即得。
\end{proof}

\begin{definition}[拐点]\label{def:g12:deriv:inflection}
$f$ 的曲线穿过其切线的点称为
\emph{拐点}\index{拐点}。对二阶可导的 $f$，
拐点即 $f''$ 变号的点。
\end{definition}

\begin{example}\label{ex:g12:deriv:cube}
$f(x) = x^3$ 的 $f''(x) = 6x$：$f$ 在 $\intoc{-\infty}{0}$ 上为凹，
在 $\intco{0}{+\infty}$ 上为凸，原点是拐点，
此处曲线穿过其切线（即 $x$ 轴）。
\end{example}

\begin{omfigure}
\begin{tikzpicture}
\begin{axis}[omaxis,
  width=7.6cm, height=5.6cm,
  xmin=-1.7, xmax=1.7, ymin=-2.9, ymax=2.9,
  xtick={-1,1}, ytick={-1,1}]
  \addplot[omProp, thick, domain=-1.2:1.2] {0};
  \addplot[omDef, thick, domain=-1.4:1.4] {x^3};
  \fill (axis cs:0,0) circle (1.5pt);
  \node[font=\small, anchor=east] at (axis cs:-0.45,-1.7) {凹};
  \node[font=\small, anchor=west] at (axis cs:0.45,1.7) {凸};
\end{axis}
\end{tikzpicture}

{\small 函数 $x \mapsto x^3$ 在原点处的拐点：曲线
穿过其切线（橙色），左侧为凹，右侧为凸。}
\end{omfigure}

\begin{example}[凸性不等式]\label{ex:g12:deriv:ineq}
指数函数在 $\R$ 上为凸函数（其二阶导数即自身，恒为正）。
它在 $0$ 处的切线为 $y = 1 + x$，于是
\[
\eu^x \geq 1 + x \quad \text{对所有 } x \in \R .
\]
此类\emph{切线不等式}是凸性的标准应用。
\end{example}

\begin{method}[利用凸性]\label{met:g12:deriv:convexity}
\begin{itemize}
  \item 证明不等式 $f(x) \geq ax + b$ 时，将该直线视为
        凸函数 $f$ 的切线，并利用
        \cref{thm:g12:deriv:convexity}(3)。
  \item 求拐点：解方程 $f''(x) = 0$，\emph{并}检验
        $f''$ 在该处是否变号。
  \item 在变化表中，凸性可细化曲线的草图
        （即曲线的弯曲方向）。
\end{itemize}
\end{method}

\section{练习}

\begin{exercise}[$\star$]\label{exo:g12:deriv:1}
求下列函数在其定义域上的导数：
\[
f(x) = x^3 - 5x + 2, \quad
g(x) = \frac{x}{x^2+1}, \quad
h(x) = (2x+1)^5, \quad
k(x) = \sqrt{x^2 + 1}.
\]
\end{exercise}

\begin{exercise}[$\star$]\label{exo:g12:deriv:2}
求曲线 $f(x) = x^2 - 3x + 1$ 在
横坐标为 $a = 2$ 的点处的切线方程，并确定曲线
相对于该切线的位置关系。
\end{exercise}

\begin{exercise}[$\star$]\label{exo:g12:deriv:3}
利用导数的定义，计算
\[
\lim_{x \to 0} \frac{(1+x)^{100} - 1}{x}
\qquad\text{和}\qquad
\lim_{h \to 0} \frac{\sqrt{4+h} - 2}{h}.
\]
\end{exercise}

\begin{exercise}[$\star\star$]\label{exo:g12:deriv:4}
研究函数 $f(x) = \dfrac{x^2 + 3}{x - 1}$ 在其定义域上的
极限、导数、变化表与局部极值。
\end{exercise}

\begin{exercise}[$\star\star$]\label{exo:g12:deriv:5}
制作一个无盖的箱子，所用材料为边长
$30\,$厘米的正方形硬纸板，在四角各剪去边长为 $x$ 的相等正方形，再将
侧面折起。求使箱子体积最大的 $x$。
\end{exercise}

\begin{exercise}[$\star\star$]\label{exo:g12:deriv:6}
设 $f(x) = x^4 - 4x^3 + 10$。
\begin{enumerate}
  \item 求 $f''$，并确定 $f$ 的凹凸区间与
        拐点。
  \item 证明：横坐标为 $2$ 的拐点处的切线
        在该点穿过曲线。
\end{enumerate}
\end{exercise}

\begin{exercise}[$\star\star$]\label{exo:g12:deriv:7}
利用切线不等式，证明对所有 $x > 0$，
\[
\sqrt{x} \leq \frac{x + 1}{2},
\]
并指出何时取等号。（提示：平方根函数是凹函数。）
\end{exercise}

\begin{exercise}[$\star\star\star$]\label{exo:g12:deriv:8}
设 $f$ 在 $\R$ 上为凸函数且可导，若 $f'$
在某点 $a$ 处为零，证明 $f(a)$ 是 $f$ 在
$\R$ 上的\emph{全局}最小值。
\end{exercise}

\begin{exercise}[$\star\star\star$]\label{exo:g12:deriv:9}
证明：在周长为 $p$ 的所有矩形中，正方形的
面积最大。再证明：在面积为 $A$ 的所有矩形中，正方形的
周长最小，并解释这两个结论
之间的联系。
\end{exercise}

\section{问题：救生员的计算}

\begin{problem}[{周末问题 --- 最快的救援路径遵循
光的斯涅尔定律，凸性可验证每个最优解，
而理想的汽水罐恰好高与直径相等}]
\label{pb:g12:deriv:1}
一名救生员看到一位遇险的游泳者 --- 他位于沿海滩斜下方
的水中。跑步比游泳快，所以
直线路径\emph{不是}最快的路线。使时间最短的路线
在岸线处折弯，其折弯规律恰是光线进入水中时发生折射所遵循的规律：
大自然也懂得求导。
本题练习链式法则（\cref{thm:g12:deriv:chain}），
完成救援，运用凸性所给出的不等式
（\cref{thm:g12:deriv:convexity}），
并设计一个铁皮罐。

\medskip
\noindent\textbf{第一部分 --- 熟练运用链式法则。}
\begin{enumerate}
  \item 求导：$\left(3x^2 + 1\right)^5$；\
        $\sqrt{x^2 + 9}$；\ $\dfrac{1}{x^2 + 1}$。
  \item 求 $y = \sqrt{x^2 + 9}$ 在 $x = 4$ 处的切线。
  \item 作出函数
        $f(x) = \dfrac{x}{x^2 + 1}$ 在 $\R$ 上的变化表（含极值）。
  \item 对 $g(x) = x^3 - 3x^2 + 4$，给出变化表\emph{与}
        凹凸性表（$g''$），并求拐点
        （见\cref{def:g12:deriv:inflection}）。
  \item 求 $g$ 在其拐点处的切线，
        并证明 $g(x) - (\text{切线}) = (x - 1)^3$：
        这一符号变化说明切线在拐点处
        如何与曲线相交？
\end{enumerate}

\medskip
\noindent\textbf{第二部分 --- 救援。}
岸线沿 $x$ 轴；救生员站在沙滩上，位置为
$A(0, 30)$，游泳者在 $B(40, -20)$ 处等候（单位：米）。
救生员在沙滩上以 $5$ 米/秒奔跑，以 $2$ 米/秒游泳，并在
她所选的点 $(x, 0)$ 处入水。
\begin{enumerate}[resume]
  \item 建模：写出救援总时间 $T(x)$，并说明为何
        只需考虑 $x \in \intcc{0}{40}$。
  \item 求 $T$ 的导数（两次用到链式法则）。
  \item 设 $\theta_1$ 为奔跑段与岸线的
        \emph{垂线}所成的角，$\theta_2$ 为游泳段的
        角。证明：
        $\sin\theta_1 = \frac{x}{\sqrt{x^2 + 900}}$，且
        条件 $T'(x) = 0$ 可写为
        \[
        \frac{\sin\theta_1}{5} = \frac{\sin\theta_2}{2}
        \]
        --- \emph{斯涅尔定律}，其中以两种速率
        代替光在空气与水中的速率。
  \item 线段 $[AB]$ 与岸线在
        $x = 24$ 处相交。求 $T(24)$、$T(40)$（沿海滩跑
        再直游过去），$T(0)$：几何上的直线是否最快；
        两种极端策略中
        是否有最快者？
  \item 对 $T'$ 用二分法定位最优入水点
        （$\approx 33.7$ 米）：$x^*$ 精确到米，
        最快用时 $T(x^*)$ 精确到十分之一秒。
        最优折弯比直线路径
        节省多少时间？
  \item 为何该驻点是\emph{全局}最小值点？
        （$T$ 的每一项都是凸函数 --- 在本阶段不加证明地承认：凸函数之和
        仍是凸函数 --- 并应用
        \cref{exo:g12:deriv:8}。）
  \item 费马原理指出：光总是沿着所需
        \emph{时间}最少的路径传播。试推导光线
        为何光线进入水中时恰在水面处发生折射
        （水中光速较慢），并指出由此可解释的一个日常现象：
        插入水杯中的铅笔看起来
        发生弯折。
\end{enumerate}

\medskip
\noindent\textbf{第三部分 --- 凸性的收获。}
\begin{enumerate}[resume]
  \item 证明 $x \mapsto x^2$ 为凸函数，并证明
        中点不等式
        $\left(\frac{a + b}{2}\right)^2 \leq
        \frac{a^2 + b^2}{2}$，并给出两种证明：一种由凸性得出，
        另一种由展开 $(a - b)^2 \geq 0$ 得出。
  \item 由此推出：\emph{平方平均数}
        $\sqrt{\frac{a^2 + b^2}{2}}$ 不小于算术
        平均数，并写出完整的不等式链 ---
        调和平均数 $\leq$ 几何平均数 $\leq$ 算术平均数 $\leq$
        平方平均数 --- 并在 $a = 2$，$b = 8$ 时验证四者。
  \item 切线位于曲线下方：$x \mapsto \frac1x$ 为凸函数，
        在 $\intoo{0}{+\infty}$ 上；写出它在 $1$ 处的切线，并
        由此推出不等式
        $\frac1x \geq 2 - x$ 对所有 $x > 0$ 成立，并指出
        何时取等号？
  \item 现实中的拐点：疫情期间，
        累计病例数呈 S 形曲线。拐点处发生了什么，
        以及为何流行病学家格外关注这一日期？
        请结合拐点两侧
        二阶导数的符号加以说明。
\end{enumerate}

\medskip
\noindent\textbf{第四部分 --- 理想的罐子。}
一个圆柱形罐子须用最少的金属容纳 $V = 330$ cm$^3$，
表面积 $S = 2\pi r^2 + 2\pi r h$，约束为
$\pi r^2 h = V$.
\begin{enumerate}[resume]
  \item 写出 $S(r) = 2\pi r^2 + \dfrac{2V}{r}$，
        求导，并计算 $V = 330$ 时的最优半径与
        高度（精确到毫米）。
  \item 证明一般性的优美结论：在最优情形下，
        $h = 2r$ --- 理想的罐子恰好高与
        直径相等。
  \item 检验 $S''(r) > 0$，并借助
        \cref{exo:g12:deriv:8}，再次得出该最优点是
        全局最小点。实际的汽水罐明显高于其直径：
        请说出一个非数学方面的原因。
  \item 结尾 --- 优化流程：建模，求
        导数，解出驻点，借助
        凸性或变化表加以判定，最后解释结果。将此流程依次用于
        救生员问题、罐子问题，以及十一年级的最强梁
        （\cref{pb:g11:deriv:1}）--- 并说明费马原理
        揭示了还有谁在运用这一流程。
\end{enumerate}
\end{problem}
